\scanpage{128}
\margin{\fbox{\parbox{1.9cm}{Show that one \unsure{can reduce} to finite dimension}}}
The first part of the theorem is summarized in the
formulas (1) to (4), \deleted{and} \inserted{besides} the assertions concerning
the vector topologies of the spaces $\mathcal L^p(\mathcal H)$.
\inserted{Let us show \deleted{\unsure{first}} first that the latter are very easy consequences \deleted{\unsure{of|of the}} \unsure{of the formulas} (3) and (4)\ldots}
Formula (1) is obvious, \unsure{for|The} \deleted{\unsure{one knows that}}
\deleted{\unsure{indeed} \unsure{in the} case $p<+\infty$, and} \deleted{\unsure{for $p<+\infty$}}, one knows (theorem
of Riesz) that $|A|$ and $|A^*|$ are \deleted{\unsure{of the same sequence}} \inserted{unitarily equivalent}
(?) \deleted{\unsure{eigenvalues}}. Formula (2) is a case
subsumed by formula (6), but \deleted{\unsure{can be|can be}} \inserted{\unsure{one can|I can}}
give a proof \deleted{\unsure{nevertheless direct}}
\deleted{prove simply in a direct way}
direct, proceeding as in $[\,]$, .

Indeed, if $(\unsure{k_i|\alpha_i})$ \deleted{\unsure{respectively $(\beta_i)$}} denotes the
sequence of \inserted{eigen}values of \deleted{$\unsure{\lvert BA\rvert|\lvert BA\rvert^2}$} $|A|$
\inserted{and $(\unsure{\beta_i|h_i})$ the sequence of eigenvalues of $\unsure{\lvert BA\rvert|\lvert B\rvert A\rvert|\lvert BA\rvert^2}$ (respectively $|AB|$)} \deleted{\unsure{respectively $|BA|$}},
arranged by magnitude
in decreasing order (each repeated as many times
as its multiplicity requires) one has
\[
\beta_i\leq\|B\|\alpha_i\quad\text{for every }i.\quad(\text{place cited, page \quad, lemma \quad}),
\]
whence at once formulas (2).

The \deleted{\unsure{last} of formulas (3) is obvious, the}
The formulas (3) \deleted{can} \inserted{can} be regarded \deleted{of formula} as
\uline{first} \deleted{\unsure{is an immediate consequence}}
\unsure{consequences|consequence} of formula (6), (which we shall prove independently) in the
following way. If $p'$ is such that $\frac{1}{p}+\frac{1}{p'}=1$,
and if $B\in\mathcal L^{p'}(\mathcal H)$, \deleted{and such that}
\inserted{$\|B\|_{p'}\leq1$}, then
one has by virtue of (6) $\|AB\|_1\leq\|A\|_p$, whence
\[
|\operatorname{Tr}(AB)|\leq\|A\|_p,\quad\text{and}\quad
\operatorname*{Sup}_{\|B\|_{p'}\leq1}|\operatorname{Tr}AB|\leq\|A\|_p
\]
